\subsubsection{Agmon's Condition and Coercive Problems}\label{subsec:agmon-condition-coercive-problems}

We consider a scalar elliptic operator of order $2m$, given in divergence form as in \eqref{eq:general-even-order-divergence-form-definition}:
\begin{equation}\label{eq:agmon-coercive-divergence-operator}
L(\bm{x},D)u=\sum_{|\alpha|,|\beta|\leq m}(-1)^{|\alpha|}D^\alpha(a_{\alpha\beta}(\bm{x})D^\beta u),
\end{equation}
where the $a_{\alpha\beta}$ are continuous on $\overline\Omega$ and the ellipticity condition
\begin{equation}\label{eq:agmon-principal-ellipticity-condition}
\sum_{|\alpha|=|\beta|=m}\bm{\xi}^\alpha a_{\alpha\beta}(\bm{x})\bm{\xi}^\beta>0
\end{equation}
holds throughout $\overline\Omega$. Moreover, we consider $p$ normal boundary-value operators $B_j(\bm{x},D)$, with coefficients of class $C^{m-m_j}(\partial\Omega)$, where $m_j<m$ is the order of $B_j$. In general, $p$ can take any value between $0$ and $m$. We define
\begin{equation}\label{eq:agmon-essential-boundary-admissible-space}
V=\{u\in H^m(\Omega)\mid B_j(\bm{x},D)u=0\text{ on }\partial\Omega,\ j=1,\ldots,p\}.
\end{equation}


We consider the quadratic form
\begin{equation}\label{eq:agmon-boundary-constrained-quadratic-form}
a(u,v)=\int_\Omega\sum_{|\alpha|\leq m,|\beta|\leq m}a_{\alpha\beta}(\bm{x})D^\beta uD^\alpha v\,\mathrm{d}\bm{x},
\end{equation}
and we ask for conditions under which this form is coercive on $V$:
\begin{equation}\label{eq:agmon-boundary-constrained-coercive-estimate}
a(u,u)\geq c_1\|u\|_m^2-c_2\|u\|_0^2\quad\forall u\in V.
\end{equation}


If the form is coercive, we can apply the Lax--Milgram lemma to conclude that, for $\lambda$ large enough, the equation
\begin{equation}\label{eq:agmon-shifted-constrained-variational-equation}
a(u,v)+\lambda(u,v)=(f,v)\quad\forall v\in V
\end{equation}
has a unique solution $u\in V$ for every $f\in V\prime$. It is then clear that $L(\bm{x},D)u+\lambda u=f$ in the sense of distributions, and that $B_j(\bm{x},D)u=0$ on the boundary. In addition $u$ will satisfy $m-p$ ``natural\textquotesingle\textquotesingle\ boundary conditions, which arise in a similar way as the Neumann boundary condition in Section~\ref{subsec:neumann-problem}.

The condition guaranteeing coercivity is known as Agmon\textquotesingle s condition. Consider a point $\bm{x}_0\in\partial\Omega$; we may orient our coordinate system in such a way that $\bm{x}_0$ is the origin and the inner normal points in the $x_n$ direction. We then consider the constant coefficient problem $L^p(0,D)u=0$ in the half-space $x_n>0$ with boundary conditions $B_j^p(0,D)u=0$ for $j=1,\ldots,p$. We shall use the notation $\bm{x}=(\bm{x}^{\prime},x_n)$, where $\bm{x}^{\prime}\in\mathbb{R}^{n-1}$, and correspondingly we write $\alpha=(\alpha^{\prime},\alpha_n)$ for a multi-index $\alpha$. We now pick any $\bm{\xi}^{\prime}\in\mathbb{R}^{n-1}\backslash\{0\}$ and consider the ODE
\begin{equation}\label{eq:agmon-principal-half-line-ode}
L^p\left(0,i\bm{\xi}^{\prime},\frac{\mathrm{d}}{\mathrm{d}t}\right)v(t)=0,\ t>0,
\end{equation}
with initial conditions
\begin{equation}\label{eq:agmon-principal-half-line-boundary-conditions}
B_j^p\left(0,i\bm{\xi}^{\prime},\frac{\mathrm{d}}{\mathrm{d}t}\right)v(0)=0,\quad j=1,\ldots,p.
\end{equation}


\begin{Definition}
\label{def:agmon-decaying-half-line-positivity}
We say that Agmon\textquotesingle s condition holds if for any $\bm{\xi}^{\prime}\in\mathbb{R}^{n-1}\backslash\{0\}$ and any nonzero solution $v$ of \eqref{eq:agmon-principal-half-line-ode} with \eqref{eq:agmon-principal-half-line-boundary-conditions} such that $v$ tends to zero exponentially as $t\to\infty$, we have the inequality
\begin{equation}\label{eq:agmon-decaying-principal-half-line-positivity}
\int_0^\infty\sum_{\substack{|\alpha^{\prime}|+k=m\\|\beta^{\prime}|+l=m}}a_{(\alpha^{\prime},k)(\beta^{\prime},l)}(0)(\bm{\xi}^{\prime})^{\alpha^{\prime}}(\bm{\xi}^{\prime})^{\beta^{\prime}}\frac{\mathrm{d}^kv(t)}{\mathrm{d}t^k}\frac{\mathrm{d}^lv(t)}{\mathrm{d}t^l}\,\mathrm{d}t>0.
\end{equation}

\end{Definition}

\begin{Remark}
\label{rem:agmon-dirichlet-vacuity}
If $p=m$ and the complementing condition holds, then Agmon\textquotesingle s condition is vacuously true. Indeed, if $p=m$, then, by Lemma~\ref{lem:dirichlet-system-tangential-triangular-reduction}, the boundary conditions are equivalent to Dirichlet conditions. In fact, Dirichlet conditions always satisfy the complementing condition; see Problem~\ref{ex:scalar-dirichlet-complementing-condition}.
\end{Remark}
